\head{第13.4段：有限点邻域、一般基底与解析拓扑}
\source{法稿物理80页中段至82页末；英编稿物理52页末至53页末；旧注D17-M081--093}

\audit{V21-REM-E01}{段落次序与原文层次}{本段五项注全部逐段对读，没有用补证替代译文}
法稿纸面依次为a)、b)、d)、c)、e)，页边另有移排箭头；英文编订稿及旧中文依编号排成a)--e)。本版沿用这一编号次序，保留每一项的完整论述，并明确记录纸面次序。法稿的手插、明确删句和不可辨覆打分别处理；英文译者的问号没有一概算作法文缺文。

本次可由法稿直接确认的内容包括：a)中的$X_{k'}$及位于原点上方的复数“点”；c)中中心$C$的余维$n$、闭点集的几何点$\Phi(\bar k)$、手插连接语及“我们已经证明的形式”；d)中几何纤维的不可约与$(R_2)$条件、$s\in S$、非光滑点集$Z$及连接维数不等式的动词；e)中的$U^{\mathrm{an}}$和上标为$n$的$\mathbb C^n$。后两项在英文中分别误作$Y^{\mathrm{en}}$、$U^{\mathrm{en}}$、$C_n$。

两项订正不能伪称为原稿本来就写对了：c)末次比较$L$的余维，法英均写$n$，须改为$n-1$；e)开头法稿字面为“初等”，正文根据13.2及任意维数的塔明确为“多重初等”，并在脚注保留原字面。此外，对$i\ne1$的同伦群表述按正次数理解成$i\geq2$，另说明连通性，不把$\pi_0$当作一个普通同伦群。

\audit{V21-REM-E02}{有限域下降与几何不可约的范围}{补足所用条件，并分开原定义与删去连通条件的弱化}
\lead{a)的有限扩张并非有限域上的原定理。}
设$k$有限，$\bar k$为其代数闭包。在一个给定的光滑仿射有限型邻域中，把所考察点换成一个闭特化$z$。$z$在$\bar k$上有有限多个几何点；将下文已补证的有限点版本应用于它们，可取一个同时包含这些点的多重初等开集。这个开集、有限长度的初等态射塔、各步的射影紧化和有限平展边界，都是有限表示的数据。写出有限多个仿射图、方程、粘合和射影嵌入后，只涉及有限多个$\bar k$中的系数；它们包含在某个有限扩张$k'/k$中。

必要时增大$k'$以同时定义这些有限数据及其等式。光滑性、有限平展性、满射性和相应开闭嵌入关系可沿忠实平坦的域扩张检验；几何纤维的连通性按几何点检验，亦在此下降。各步有限平展边界的度数在连通底空间上恒定。因此得到定义在$k'$上的初等塔，而不只是一个方程形式相似的候选塔。所有$z$上方的点都属于其起始开集，这也可以在$\bar k$上检验。

对原来的非闭点$x$，有限态射$X_{k'}\to X$的上升性质保证：每个$x$上方的点都有一个位于$z$上方的特化。含全部这些特化的开集也含它们的全部泛化，故仍同时包含$x$上方的点。这个论证说明a)的有限扩张范围；没有消去$k'$，没有证明有限域上可直接令$k'=k$。正文保留作者关于去掉无限性假设的推测身份。

\lead{原定义的几何不可约性。}
对任意代数闭扩域，将一个非空初等塔基变换过去。每一步的紧化纤维是光滑连通曲线，删去非空有限点集后仍为连通非空的光滑曲线。初等态射是光滑满射，故为开映射。若底空间连通而总空间分成两个非空开闭部分，连通纤维必完全落在其中一部分；两部分的像便给出底空间的两个互不相交的非空开集，矛盾。由点开始归纳，各层光滑且连通。域上有限型的光滑概形正规，有限多个不可约分量互不相交并开闭，所以连通即不可约。这证明原定义下多重初等概形的几何不可约性。零次复合为恒等态射，其定义域是$\Spec k$。

\lead{仅删除“连通”仍有零维边界。}
13.2没有正维假设，13.1却要求每一步的几何纤维纯一维。设$L/k$为非平凡的有限可分扩张，例如$k=\mathbb R$、$L=\mathbb C$，取$X=\Spec L$。它在$k$上光滑，唯一点只有$X$自身这一非空开邻域。删去连通条件之后，每一个非空初等态射仍是相对维数1的光滑满射；长度$r$的塔在$k$上的相对维数为$r$。因而到达零维$X$只能取$r=0$，但零次复合只给$\Spec k$，不等于$\Spec L$。

所以原b)中“只删连通即可使弱化结论成立”的无条件断言，在零维时失败。这是对表述范围的反例，不是对13.2原有几何不可约版本的反例。可以考虑把末尾的有限平展态射也列入允许的塔，或单独规定零维起始对象；但这属于进一步修改定义。本版明确指出所需补约，不假称仅删一个词已经证明所有维数的推广，也不据此声称某一个新定义的全部高维结论已经核定。

\audit{V21-REM-E03}{有限个光滑点的共同多重初等邻域}{在无限完美域上给出完整归纳补证，兼顾非有理点}
以下证明为本轮译者补证。其几何连通性输入单独注明；不是宣称原稿8.7(c)、8.8的缺续已经恢复。

\lead{精确命题。}
设$k$为无限完美域，$X$为几何不可约的$k$-预概形，$\Phi$为其光滑点组成的非空有限集合，并包含在某个仿射开集中。对任意包含$\Phi$的开集$W$，存在开集$U$，使
\[
\Phi\subset U\subset W,\qquad U\longrightarrow\Spec k
\quad\hbox{为多重初等态射。}
\]
“基本邻域系”正是这个对每个$W$都成立的结论。这里明确排除空的$\Phi$：若取$\Phi=W=\varnothing$，就不能在$W$内找到到$\Spec k$的非空多重初等塔。原13.4(c)没有明说这一边界，本版正文已经补明；上述非空范围足以覆盖13.2及这里的共同邻域构造。

\begin{analysisproof}
\lead{第一步：有限型约化和归纳起点。}
在包含$\Phi$的仿射开集$A$内，取$W\cap A$与光滑轨迹的交。其闭补集由某个理想$I$定义，而$I$不包含于任何$\Phi$中点的素理想。有限素理想回避给出$f\in I$，在$\Phi$的每一点均不为零。因此
\[
\Phi\subset D(f)\subset W\cap A
\]
且$D(f)$光滑。这个主开集拟紧，光滑性给出局部有限表示，故$D(f)$为$k$上有限型的仿射概形。对每个原点，在其闭包与$D(f)$的交中取一个闭点；这就是一个仍在$D(f)$内的闭特化。含这些有限个闭点的开集也含原来的点。故可将$X$换成一个光滑几何整的仿射有限型概形，将$\Phi$换成有限闭点集。令$n=\dim X$，并将$\Phi$赋予约化结构。因为$k$完美，$\Phi\to\Spec k$有限平展，因而$\Phi(\bar k)$确是一个有限集合。

以下按$n$对这个有限点命题归纳，不能只假设单点13.2后直接把多个邻域相交而认定其仍为多重初等。若$n=0$，光滑几何整性使$X=\Spec k$。若$n=1$，取$X$的正规射影完备化$\overline X$。完美域上的正规曲线光滑；其几何整性由原来的稠密几何整开集保证。约化边界$D=\overline X\setminus X$有限平展，且非空：正维的仿射概形不能同时为域上的固有整概形。因此$X\to\Spec k$本身为初等态射。下面设$n\geq2$。

\lead{第二步：选择射影环境和坏闭集。}
取$X$的射影闭包再正规化。有限型域上正规化有限，故所得$\overline X$仍为射影概形；在光滑开集$X$上正规化是同构。它几何整且正规。取充分高次的超曲面截面，包含$\overline X\setminus X$而避开$\Phi$，并令$Y$为该截面的约化闭子概形。这样的截面存在：闭补集的理想层充分扭转后，可以在有限闭子概形$\Phi$上指定处处非零的值。$Y$纯余维1；由于$k$完美，$Y$的一般点光滑。令
\[
Z=\operatorname{Sing}(\overline X)\cup\operatorname{Sing}(Y).
\]
正规性及$Y$的约化纯维性给出$\dim Z\leq n-2$，并且$Z\subset Y$。特别地，$\Phi\cap Y=\varnothing$。重新取一个射影嵌入后，再作次数$d\geq2$的Veronese嵌入，记新的环境为$\mathbf P^N$。

\lead{第三步：各个几何点的良好线性截面。}
对$x_i\in\Phi(\bar k)$，考虑经过$x_i$、余维为$n-1$的线性子空间$L$。本版V21-BER-E03的一阶截面补证，分别应用于$\overline X\setminus Z$和$Y\setminus Z$，给出非空开参数条件，使其截面分别光滑、纯维数1及光滑、维数0。对后一概形，$x_i$根本不在其上，因此没有固定点处的低维分量障碍。

这些一般$L$还可以避开$Z$。具体地，对$z\in Z$，经过$x_i,z$的余维$n-1$子空间，在所有经过$x_i$的参数空间中具有余维$n-1$；而$\dim Z\leq n-2$。对应关联集的像因$Z$射影而闭，且维数小于整个参数空间。因此避开$Z$是另一个非空开条件。两类条件相交后，得到非空开集$O_i$，其中每个$L$都具有13.3所需的三项性质。

\lead{第四步：同一个中心同时适合所有几何点。}
令$G$参数化$\mathbf P^N$中余维$n$的线性中心$C$。与$Y\cup\Phi$不交是其中一个非空开条件。先暂不施加这一条件：对固定$x_i$，在仅要求$C$不含$x_i$的最大定义域上有态射
\[
\sigma_i\colon C\longmapsto L_i=\langle C,x_i\rangle.
\]
它是到所有经过$x_i$、余维$n-1$子空间的满射：给定$L_i$，取其中不含$x_i$的任意超平面作为$C$即可。更准确地说，若$x_i$对应向量直线$\ell_i$、$L_i$对应向量空间$V_i$，则纤维为$\ell_i$在$V_i$中的补空间全体。局部选择一个补空间后，这一选择空间与维数$N-n+1$的仿射空间
\[
\operatorname{Hom}(V_i/\ell_i,\ell_i).
\]
相识别；这种识别依赖所选的补空间。相对地，$\sigma_i$是上述Hom丛的仿射丛，因而光滑且满射，$\sigma_i^{-1}(O_i)$是非空开集。然后才把它与中心避开$Y\cup\Phi$的非空开集相交；不可约性保证交仍非空。没有声称加上所有避开条件之后，$\sigma_i$仍满射到全部$L_i$。

还须另加一个独立条件：
\[
T=C\cap\overline X\quad\hbox{非空且有限平展于 }k.
\]
不能从各$L_i$截面良好就自动推出它。用完全线性系的一阶截面满射，按V21-BER-E03的值及微分秩计算而不设基点，一般的余维$n$截面在$\overline X\setminus Z$上光滑、维数0；一般中心避开$Z$。射影维数定理又保证$C\cap\overline X$非空。对应的泛关联族射影，去掉非平展及正维纤维的闭像后，就得到所需的非空开参数条件。

这些条件可同时在$k$上定义。令$G_0$为中心避开$\Phi$的开集，在$S=G_0\times_k\Phi$上有统一的张成空间态射。先以射影关联族$Z\cap L$的闭像排除遇到$Z$的参数。在其补开集内，$\overline X$及$Y$的关联族都处于光滑环境中；用相对Jacobian秩条件，以及相对维数分别为1和0的条件，定义良好轨迹。关联族射影，坏轨迹的像闭，因而得到$S$中的一个在$k$上定义的好开集$O$。在每一个几何$x_i$对应的纤维上，前述固定点补证保证其非空；在排除$Z$的开集中，满微分秩也直接保证正确的相对维数。

于是$S\setminus O$为坏闭集。投影$S\to G_0$有限，其像闭；像的补集恰要求每一个几何点$x_i$均满足良好条件。将这个开集与中心避开$Y$、以及$T$非空有限平展的开集相交。几何基变换后，它是不可约格拉斯曼簇中有限个非空开集的交，故非空；于是得到定义在$k$上的非空开集$G^\circ$。这一步利用统一的关联族定义条件，不是先任取各个$O_i$之后，未经说明就假定它们能下降。

格拉斯曼簇有定义在$k$上的稠密标准仿射开胞腔。无限域上仿射空间的$k$-有理点稠密，因此$G^\circ(k)\ne\varnothing$。选择其中一个$C$。此时$C$在$k$上定义，但各$L_i$只须在相应几何点的域上定义；证明从未把非有理点擅自当作$k$-有理点。

\lead{第五步：共同的初等曲线族。}
令$Q\simeq\mathbf P^{n-1}_k$参数化包含$C$、余维$n-1$的子空间，取相应的射影关联族
\[
p\colon\mathcal X=\{(z,L):z\in\overline X\cap L\}\longrightarrow Q.
\]
在$\overline X\setminus C$上，它由以$C$为中心的线性投影给出，关联族的另一投影是同构。每个$\Phi$中的闭点投影为$Q$中的闭点；记这些点的有限集合为$\Psi$。其几何纤维正是已选好的$\overline X\cap L_i$。这些纤维避开$Z$，其定义方程的相对微分独立，故Jacobian判据使$p$在它们上光滑、相对维数1。坏轨迹的射影像闭，因此可以在$Q$中取一个同时含$\Psi$的开集$B$，使$p_B$光滑。再删去相对维数不等于1的开闭层的射影像，即可使$p_B$处处相对维数1，仍不删去$\Psi$。

这里还需要\emph{几何连通性}，不能从光滑性直接推出。所用的连通Bertini定理是：代数闭域上，若$V$为不可约射影簇、$f\colon V\to\mathbf P^N$为态射，则在
\[
\dim L\geq N+1-\dim f(V)
\]
时，$f^{-1}(L)$连通。将它应用于$\overline X$的嵌入和每个余维$n-1$的$L$，即得每个几何曲线纤维连通。参见Martinelli--Naranjo--Pirola，\emph{Connectedness Bertini Theorem via numerical equivalence}，定理1.1；作者稿\url{https://arxiv.org/abs/1412.1978}。这是本补证明确引用的外部定理，不是本版重新证明的一个结论。

同理，$Y$的关联子概形在$\Psi$上给出光滑零维截面，故缩小$B$后在$B$上有限平展。另一边界为$T\times_k B$，也有限平展且满射；它与$Y$的关联子概形不交，因为$T\cap Y=\varnothing$。$B$是射影空间的非空开集，故几何不可约；边界的有限平展度数在其上恒定。删去这两个闭子概形，得到
\[
U_0\longrightarrow B
\]
这一初等态射；其紧化为$\mathcal X_B$。它的每个纤维是连通光滑射影曲线删去有限多个点，因而非空。通过线性投影识别，$U_0$是原来$X$中的开集，并含$\Phi$。非空边界由$T\times B$保证；不能只指望$Y$截面在每个纤维上非空。

\lead{第六步：在基底上用有限点归纳。}
在$B\subset\mathbf P^{n-1}$中，有限闭点集$\Psi$有共同仿射开邻域$B'$。例如可取包含$\mathbf P^{n-1}\setminus B$而避开$\Psi$的充分高次超曲面，其补集既仿射又包含于$B$。对维数$n-1$的光滑几何整概形$B'$及其有限点集$\Psi$应用归纳假设，得到多重初等开集$V\subset B'$，且$\Psi\subset V$。令
\[
U=U_0\times_B V.
\]
这是$X$中含$\Phi$的开集；$U\to V$因基变换仍为初等态射，而$V\to\Spec k$为多重初等态射，所以$U\to\Spec k$也为多重初等态射。最初已将工作限制在给定$W$内，故$U\subset W$。归纳完成。
\end{analysisproof}

这个补证同时解决了13.2证明中先写一般光滑点、后改用$x\in P(k)$所留下的入口问题：原点先取闭特化，再同时处理其全部几何点；无需假定原来的点是$k$-有理点。它也说明不能直接把不同点的多重初等邻域取交而省略共同构造。

\audit{V21-REM-E04}{一般基预概形的推广}{文字条件已经恢复，一般数学断言仍未核定}
法稿d)明确将$(R_2)$施于几何纤维，并明确说“不作证明”。$(R_2)$表示所有余维至多2的点处局部环正则；在此处不得误写成仅有余维1的正则性，亦不得把几何纤维条件弱化为原剩余域上的条件。$S'$在$S$上有限、纤维在相关点处光滑、$Y$避开$S'$、非光滑点集合为$Z$以及维数不等式，均逐项保留。

但是，恢复这些词句不等于证明任意基底上的结论。正文只要求在$s$的Zariski开邻域上构造整个多重初等塔；先在有限扩域或平展邻域上找到投影中心，并不能自动下降成所要求的Zariski局部结论。有限剩余域、任意基底上有限表示的使用，以及7.1中允许的基变换与这里结论是否一致，仍须单独处理。本版没有把这些未处理条件删掉，也没有凭尚未找到证明就断言已得到一个反例。该项维持“原作者提出、未附证明；本轮未完成一般情形核查”的状态。


% 本片段直接拼接进 reports/v21_section13_remarks_audit.tex。
% 只用当前项目已有宏；无导言区，无外部输入；以下证明属于本次译校补证。
\audit{V21-REM-E05}{普遍覆盖的欧氏空间结论}{原定义下成立；补充证明不属于原稿所附文字}
\source{法文物理第74页定义13.1及第81页末至第82页备注13.4(e)；英文第49、53页}

按定义13.1的条件，复数域上的$n$维多重初等概形$U$的普遍覆盖同胚于$\mathbb C^n$这一结论成立。下面的证明还给出实光滑意义的微分同胚
\[
\widetilde{U^{\mathrm{an}}}\cong_{\mathrm{diff}}\mathbb R^{2n}.
\]
原稿在此只陈述结论，以下是本次译校补证。不能仅由高阶同伦群消失，或普遍覆盖可缩，就省略这个较强的欧氏空间结论的论证。

\par\medskip\noindent\textbf{一、原定义所提供的条件。}
记初等态射塔为
\[
U=U_n\longrightarrow U_{n-1}\longrightarrow\cdots
\longrightarrow U_0=\Spec\mathbb C.
\]
每一步写成
\[
U_j=\overline U_j\setminus D_j,\qquad
\overline f_j\colon\overline U_j\longrightarrow U_{j-1},
\]
其中$\overline f_j$光滑、射影，几何纤维为连通的非空光滑曲线，而$D_j\to U_{j-1}$有限、平展且满射。原稿虽未在此另写“有限”，但$D_j$闭嵌入射影族，故它到基底的态射固有；又因平展而拟有限，故为有限态射。满射条件保证每个纤维的去点数$r_j\geq1$。同类定义也见下列参考材料[D]的第III.4.2段，其中明确要求边界覆盖没有空纤维。

从$\Spec\mathbb C$开始归纳，各$U_j$及紧化均为$\mathbb C$上的光滑、分离、有限型概形，其解析化是Hausdorff、第二可数且仿紧的光滑流形。每步相对复维数为1，故$\dim_{\mathbb C}U_j=j$。连通光滑曲线删去有限点后仍非空且连通；结合下面证明的局部平凡性，各$U_j^{\mathrm{an}}$归纳为非空、连通并道路连通。

\par\medskip\noindent\textbf{二、保持被删边界的水平分裂。}
考虑某一步的解析紧化
\[
p\colon M=\overline U_j^{\mathrm{an}}\longrightarrow
B=U_{j-1}^{\mathrm{an}},\qquad D=D_j^{\mathrm{an}}\subset M.
\]
$p$是固有的满浸没，$D$是闭子流形，且$p|_D$是有限覆盖和局部微分同胚。普通Ehresmann定理适用于$p$，参见[E]；但是，仅引用这一定理还不足以保证局部平凡化同时保持$D$。这个附加性质可直接构造如下。

令$V=\ker(dp)$，先取光滑向量丛正合列
\[
0\longrightarrow V\longrightarrow TM
\overset{dp}{\longrightarrow}p^*TB\longrightarrow0
\]
的任一光滑分裂$s_0$。在$D$上规定分裂
\[
s_D=\bigl(d(p|_D)\bigr)^{-1}
\colon p^*TB|_D\longrightarrow TD\subset TM|_D.
\]
差$s_D-s_0|_D$是
$\operatorname{Hom}(p^*TB,V)|_D$的光滑截面。由于$D$是闭子流形，借助管状邻域及单位分解可将这个截面延拓至整个$M$。将延拓加到$s_0$上，得到仍满足$dp\circ s=1$的光滑分裂$s$，且$s|_D=s_D$。于是水平分布$H=\operatorname{im}(s)$满足
\[
H|_D=TD.
\]
不能在这一步只任取一个黎曼度量而不检查所得水平分布是否保持$D$。

给定定义于紧区间的分段光滑路径$\gamma\colon[0,1]\to B$，其水平提升满足光滑常微分方程，并始终位于紧集
\[
p^{-1}\bigl(\gamma([0,1])\bigr)
\]
之内；这里使用的是紧化$p$的固有性。因此，常微分方程的延拓定理保证提升存在于整个区间。因$H|_D=TD$，从$D$开始的提升始终留在$D$中；逆向路径的平行移动也保持$D$，故正、逆平行移动同时保持其补集$M\setminus D$。这里“边界”指被删去的闭子流形，不是将$M$改成了带实边界的流形。

在基底的坐标球$W$中选择径向路径，就得到一对空间的光滑局部平凡化
\[
\bigl(p^{-1}(W),D\cap p^{-1}(W)\bigr)
\cong_{\mathrm{diff}}W\times(M_b,D_b).
\]
从而开族$U_j^{\mathrm{an}}\to U_{j-1}^{\mathrm{an}}$也是光滑纤维丛，纤维为$F_j=M_b\setminus D_b$。此论证不要求这个开族自身固有；所需的固有性来自紧化，不能直接把普通固有浸没定理无条件用于开族。

\par\medskip\noindent\textbf{三、曲线纤维的基本群与普遍覆盖。}
$M_b$是亏格$g_j$的紧连通Riemann曲面，而$D_b$有$r_j\geq1$个点。去点曲面的多边形表示表明，$\pi_1(F_j)$有生成元
$a_1,b_1,\ldots,a_{g_j},b_{g_j},c_1,\ldots,c_{r_j}$，其唯一关系为
\[
\prod_{\nu=1}^{g_j}[a_\nu,b_\nu]c_1\cdots c_{r_j}=1.
\]
因$r_j\geq1$，可以利用唯一的关系消去$c_{r_j}$，故这是秩$2g_j+r_j-1$的自由群。该秩允许等于零：$(g_j,r_j)=(0,1)$时，$F_j$是复平面，基本群为平凡群。

普遍覆盖$\widetilde F_j$带有提升的Riemann曲面结构。由单连通Riemann曲面的均匀化定理[U]，它为Riemann球面、复平面或单位圆盘之一。第一种不可能，因为球面是紧的，它的连续像不能是非紧的$F_j$。余下两种作为实光滑流形都微分同胚于$\mathbb R^2$。因此
\[
\widetilde F_j\cong_{\mathrm{diff}}\mathbb R^2,\qquad
\pi_q(F_j)=0\quad(q\geq2).
\]
这不要求$\widetilde F_j$在复解析意义下总是复平面。

\par\medskip\noindent\textbf{四、欧氏普遍覆盖的归纳构造。}
归纳假设$B=U_{j-1}^{\mathrm{an}}$的普遍覆盖
$q\colon\widetilde B\to B$满足
\[
\widetilde B\cong_{\mathrm{diff}}\mathbb R^{2j-2}.
\]
把紧化与边界一起沿$q$拉回，记为
\[
\widetilde M=M\times_B\widetilde B,\qquad
\widetilde D=D\times_B\widetilde B,\qquad
\widetilde E=(M\setminus D)\times_B\widetilde B.
\]
拉回后的紧化仍为固有的光滑满浸没。选定归纳所得的欧氏坐标，并在$\widetilde M$上取上述保持$\widetilde D$的水平分裂。对每个$b\in\mathbb R^{2j-2}$，沿规定的径向路径$t\mapsto tb$从$0$平行移动到$b$。由完整提升和常微分方程的解对初值、参数的光滑依赖，得到全局微分同胚
\[
\begin{aligned}
(\widetilde M,\widetilde D)
&\cong_{\mathrm{diff}}\mathbb R^{2j-2}\times(M_0,D_0),\\
\widetilde E&\cong_{\mathrm{diff}}\mathbb R^{2j-2}\times F_j.
\end{aligned}
\]
逆映射由相同径向路径的反向平行移动给出。这里没有假定联络平坦，也没有声称平行移动与路径无关；每个端点只使用已经规定的一条径向路径。

考察复合
\[
\mathbb R^{2j-2}\times\widetilde F_j
\longrightarrow\widetilde E\longrightarrow U_j^{\mathrm{an}}.
\]
第一箭头是乘积覆盖，第二箭头是基底覆盖$q$的拉回覆盖。为了说明复合仍为覆盖，在$U_j^{\mathrm{an}}$中取一个包含于第二箭头的均匀覆盖邻域内的单连通坐标球。其每一张提升也是单连通的球，第一箭头在每张提升上均平凡，故复合均匀覆盖原坐标球。这个说明使用了流形的局部单连通性，避免把一般空间中覆盖的复合性质当作无条件成立。

复合的定义域连通、单连通，且
\[
\mathbb R^{2j-2}\times\widetilde F_j
\cong_{\mathrm{diff}}\mathbb R^{2j}.
\]
因此它确为$U_j^{\mathrm{an}}$的普遍覆盖。归纳从$j=0$的一点开始，遂有
\[
\widetilde{U^{\mathrm{an}}}
\cong_{\mathrm{diff}}\mathbb R^{2n}
\cong_{\mathrm{homeo}}\mathbb C^n.
\]
这一步直接证明了欧氏空间结论，而非只证明可缩性。

\par\medskip\noindent\textbf{五、同伦群、自由群扩张与分类空间。}
覆盖在$q\geq2$的同伦群上诱导同构，因而
\[
\pi_q(U_j^{\mathrm{an}})=0\qquad(q\geq2).
\]
也可以用纤维丛的同伦正合列归纳证明这一点。低阶正合列为
\[
\begin{split}
\pi_2(U_{j-1}^{\mathrm{an}})
&\longrightarrow\pi_1(F_j)
\longrightarrow\pi_1(U_j^{\mathrm{an}})\\
&\longrightarrow\pi_1(U_{j-1}^{\mathrm{an}})
\longrightarrow\pi_0(F_j).
\end{split}
\]
左端群为零，而右端为单点，所以得到短正合列
\[
1\longrightarrow F_{\,2g_j+r_j-1}
\longrightarrow\pi_1(U_j^{\mathrm{an}})
\longrightarrow\pi_1(U_{j-1}^{\mathrm{an}})
\longrightarrow1,
\]
其中$F_m$表示秩$m$的自由群。纤维基本群的单射性使用了基底的$\pi_2=0$，不能仅从纤维连通推出。所用同伦正合列及纤维丛的同伦提升性质见[H]。

由此总基本群是有限次自由群扩张。这里并未证明、也不需要声称这些扩张全部分裂。光滑流形具有CW同伦型，结合连通性和可缩的普遍覆盖，$U^{\mathrm{an}}$成为
$K(\pi_1(U^{\mathrm{an}}),1)$。把$\pi_1(U^{\mathrm{an}})$视为离散拓扑群，甲板变换作用使普遍覆盖成为通常同伦分类意义的普遍主$\pi_1$-丛；这不是宣称分类映射在不计同伦时严格唯一。

\par\medskip\noindent\textbf{六、适用边界及必要限制。}
\begin{enumerate}[label=\arabic*.]
\item 非空边界不能省去。若允许$r=0$，则
$\mathbb P^1_{\mathbb C}\to\Spec\mathbb C$
满足其余光滑、射影、连通曲线条件，但其普遍覆盖仍为球面，
$\pi_2\cong\mathbb Z$。原定义的满射边界正好排除了这个反例。
\item 原结论是同胚，不是双全纯同构。上面得到的加强是实光滑微分同胚。例如
$\mathbb P^1\setminus\{0,1,\infty\}$
的基本群为非交换自由群$F_2$，其复解析普遍覆盖不可能为$\mathbb C$：复平面的全纯自同构是仿射变换，而无不动点的非恒等仿射变换只能是平移；这样甲板群将交换，与$F_2$矛盾。结合[U]的均匀化分类，其普遍覆盖为圆盘。
\item 当$n=0$时，取空态射族的复合，$U\simeq\Spec\mathbb C$，解析空间与普遍覆盖都是一点，即
$\mathbb C^0=\mathbb R^0$。这正是归纳起点，不能把空概形混作这里有基本群和普遍覆盖的对象。
\item 本证明使用13.1的几何连通纤维条件。若使用13.4(b)中去掉连通条件的弱化定义，必须另查每个曲线分支是否与边界相交，并按连通分支表述基本群及普遍覆盖，不能直接照搬本证明。
\item 原稿的$i\ne1$是在正整数编号的同伦群范围内使用。严谨的高阶消失式为$i\geq2$；$\pi_0$通常是分支集合，连通性应另行说明。原稿的欧氏普遍覆盖论断在上述条件下正确，不能把原文没有写证明本身当成已发现的定理反例。
\end{enumerate}

\par\medskip\noindent\textbf{所用一手参考材料。}
下列材料的指定条目已实际打开读取。保持边界的水平分裂及其径向平凡化是以上补证所写，不能说成普通Ehresmann定理的引文已经包含全部论证。
\begin{description}[style=nextline]
\item[{[E]}] Daniel Monclair，课程讲义第1章\emph{Fibre bundles}，
Theorem 1.1.7，印刷第8页：固有满浸没为光滑纤维丛。
\par\href{https://www.imo.universite-paris-saclay.fr/~daniel.monclair/chap2.pdf}{作者课程讲义PDF（Paris-Saclay）}。
\item[{[U]}] Jacob Lurie，\emph{Surfaces and Complex Analysis (Lecture 34)}，
2009年5月3日，第1页Theorem 1、Theorem 2：单连通Riemann曲面均匀化及复平面全纯自同构。
\par\href{https://people.math.harvard.edu/~lurie/937notes/937Lecture34.pdf}{作者Lecture 34讲义PDF（Harvard）}。
\item[{[H]}] Allen Hatcher，\emph{Algebraic Topology}第4章，
Theorem 4.41（印刷第376页）、Proposition 4.48（第379页），
另见Example 4.49（第380页）：纤维丛的同伦提升、同伦正合列和覆盖的高阶同伦群。
\par\url{https://pi.math.cornell.edu/~hatcher/AT/ATch4.pdf}
\item[{[D]}] P.~Deligne讲授，J.~F.~Boutot记录，
\emph{Étale cohomology: starting points}，第III.4.2段，印刷第21--22页，
J.~S.~Milne网站提供的署名英文转译：带无空纤维有限平展边界的初等纤维化及其复合。
\par\url{https://www.jmilne.org/math/Documents/DeligneArcata.pdf}
\end{description}
